% !TEX program = lualatex
% !TeX encoding = UTF-8
\PassOptionsToPackage{unicode}{hyperref}
\PassOptionsToPackage{bookmarksnumbered,bookmarksopen}{hyperref}
\documentclass[12pt,a4paper]{article}

% ============================================================
% 한국어 설정 – Windows 시스템 글꼴 (Noto Sans KR / Noto Serif KR)
% ============================================================
\usepackage{luatexja}
\usepackage{luatexja-fontspec}

% 한국어 고딕체 (sans) – Noto Sans KR
\setsansjfont{Noto Sans KR}[
UprightFont = * Regular,
BoldFont = * Bold,
Script = Hangul,
Language = Korean
]

% 한국어 명조체 (serif) – Noto Serif KR
\IfFontExistsTF{Noto Serif KR}{
	\setmainjfont{Noto Serif KR}[
	UprightFont = * Regular,
	BoldFont = * Bold,
	Script = Hangul,
	Language = Korean
	]
}{
	% fallback: Noto Sans KR
	\setmainjfont{Noto Sans KR}[
	UprightFont = * Regular,
	BoldFont = * Bold,
	Script = Hangul,
	Language = Korean
	]
}

% 고정폭 – 영문은 Consolas, 한글은 Noto Sans KR
\setmonojfont{Noto Sans KR}[
UprightFont = * Regular,
BoldFont = * Bold,
Script = Hangul,
Language = Korean
]

% 라틴 문자 글꼴 (영문)
\setmainfont{Liberation Serif}[Ligatures=TeX]
\setsansfont{Arial}[Ligatures=TeX]
\setmonofont{Consolas}[
WordSpace={1,0,0},
HyphenChar=None,
PunctuationSpace=WordSpace
]

% ============================================================
% 기타 패키지
% ============================================================
\usepackage{amsmath,amssymb,amsfonts,mathtools}
\usepackage{geometry}
\geometry{left=2.5cm,right=2.5cm,top=2.5cm,bottom=2.5cm}
\usepackage{booktabs}
\usepackage{tabularx}
\usepackage{array}
\usepackage{hyperref}
\usepackage{url}
\usepackage{graphicx}
\usepackage{caption}
\usepackage{subcaption}
\usepackage{float}
\usepackage{siunitx}
\usepackage{bm}
\usepackage{microtype}
\usepackage{tcolorbox}

\hypersetup{
	colorlinks=true,
	linkcolor=blue,
	citecolor=blue,
	urlcolor=blue,
}

% ============================================================
% YUCT 매크로 (변경 없음)
% ============================================================
\newcommand{\Keff}{K_{\mathrm{eff}}}
\newcommand{\kappac}{\kappa_c}
\newcommand{\eps}{\varepsilon}
\newcommand{\betaf}{\beta}
\newcommand{\df}{d_{\!f}}
\newcommand{\Sodd}{S_{\mathrm{odd}}}
\newcommand{\Seven}{S_{\mathrm{even}}}
\newcommand{\Mpl}{M_{\mathrm{Pl}}}
\newcommand{\Lpl}{l_{\mathrm{Pl}}}
\newcommand{\Mk}{M_{k}}

% ============================================================
% 제목 및 저자 정보
% ============================================================
\title{\textbf{부록 J (개정판): YUCT 스케일링 양자로부터의 페르미온 질량의 반정수 양자화}\\
	\vspace{0.5cm}
	\large 현상론에서 기본 예측으로}
\author{Alexey V. Yakushev\\
	\vspace{0.5cm}
	Yakushev Research, \url{https://yuct.org/}\\
	\url{https://doi.org/10.5281/zenodo.18444598}}
\date{2026년 4월 (버전 3.0)}

\begin{document}
	
	\maketitle
	\vspace{0.5cm}
	\noindent\textbf{초록}
	
	\noindent
	본 논문에서는 야쿠셰프 통일 조정 이론 (YUCT)의 틀 안에서 표준 모형 페르미온 질량의 정밀 분석을 제시한다. 보편 스케일링 지수 \(\beta = 2/3\)를 \(2 \times (1/3)\)로 분해함으로써 기본 스케일링 양자 \(q = (3/2)^{1/3} \approx 1.1447\)을 확인한다. 이 양자는 로그 질량 스펙트럼을 지배한다. 9개의 하전 페르미온 질량은 모두 \(m_f = m_e \cdot q^{N_f}\)의 형태를 따르며, \(N_f\)는 \textbf{반정수} 값을 취한다. 이로 인해 기존 YUCT 매개변수화의 18개 자유 매개변수가 10개로 감소하고, 동시에 최대 편차가 6\%에서 1\% 미만으로 개선된다. 이러한 패턴이 우연히 나타날 확률은 \(10^{-15}\) 미만이다. 인자 \(1/3\)의 기하학적 기원을 3차원 공간에서, 반정수 단계를 스핀 통계로부터 논의한다. 개정된 양자화는 중성미자 질량에 대한 검증 가능한 예측을 제공하며, 4세대의 순차적 존재를 강력히 배제한다. 본 부록은 이전 YUCT 버전의 경험적 플레이버 분석을 대체하며, 제일원리 유도를 위한 견고한 현상론적 기초를 확립한다.
	
	\vspace{0.5cm}
	\newpage
	
	\section{서론}
	
	원래 부록 J에서는 표준 모형의 플레이버 매개변수가 현상론적 공식
	\begin{equation}
		m_f = M_{\mathrm{Pl}} \left( \frac{2}{3} \right)^{96 + k_f} \xi_f,
		\label{eq:old}
	\end{equation}
	으로 기술되었다. 여기서 \(k_f \in \{4,6,7,8,9,11,12\}\)는 정수 오프셋이고 \(\xi_f\)는 데이터에서 추출된 공명 인자이다. 수 퍼센트의 정확도이지만, 이 매개변수화는 18개의 독립적인 수치를 사용하며 근본적인 양자화 규칙에 대한 통찰을 거의 제공하지 않는다.
	
	기본 오차 스케일링 지수 \(\beta = 2/3\)가 차원 인자 \(1/3\)에서 비롯된다는 발견 (부록 Y 참조)은 로그 질량 척도상의 진정한 기본 단계가 \(\ln(3/2)\)가 아니라 \(\frac{1}{3}\ln(3/2)\)임을 시사한다. 이는 \textbf{스케일링 양자}
	\begin{equation}
		q \equiv (3/2)^{1/3} \approx 1.1447
	\end{equation}
	로 이어진다. 본 개정 부록에서는 모든 하전 페르미온 질량이 \(\ln q\)의 반정수 단위로 양자화되며, 훨씬 적은 매개변수로 전례 없는 정밀도를 달성함을 보여준다.
	
	\section{반정수 양자화의 유도}
	
	\subsection{\(\beta = 2/3\)에서 스케일링 양자로}
	
	YUCT 대수 상수는 \(\Sodd = 6/5\), \(\Seven = 4/5\)를 만족하며,
	\[
	\frac{\Sodd}{\Seven} = \frac{3}{2} = \frac{1}{\beta}
	\]
	이다. 비율 \(3/2\)는 세대 간 질량 인자를 결정한다. 그러나 보편 스케일링 법칙에서 인자 \(2/3 = (1/3) \times 2\)의 존재는 기본 기하학이 각 조정 층을 세 개의 직교하는 하위 층으로 분할함을 나타낸다. 따라서 조정 깊이 \(L\)의 한 단위는 질량을 \(q^3 = 3/2\)만큼 변화시키지만, 더 작은 단계도 허용된다.
	
	이에 따라 하전 페르미온의 질량이
	\begin{equation}
		m_f = m_e \cdot q^{N_f} \cdot \eta_f,
		\label{eq:new}
	\end{equation}
	로 주어지는 양자수 \(N_f\)를 정의한다. 여기서 \(m_e = 0.51099895\) MeV는 전자 질량이고, \(\eta_f \approx 1\)은 작은 보정 (일반적으로 \(\pm 2\%\) 이내)이다. 전자 자체는 \(N_e = 0\)의 기준점 역할을 한다.
	
	\subsection{반정수 양자화}
	
	입자 데이터 그룹 \cite{PDG2024}의 실험 질량을 사용하여 각 페르미온에 대한 최적의 \(N_f\)를 계산한다. 놀랍게도 모든 값이 \textbf{정수 또는 반정수}에 매우 가깝다. 결과는 표 \ref{tab:new_masses}에 요약되어 있다.
	
	\begin{table}[H]
		\centering
		\caption{하전 페르미온 질량의 반정수 양자화}
		\label{tab:new_masses}
		\small
		\begin{tabular}{l c c c c c}
			\toprule
			페르미온 & 실험 질량 (GeV) & \(N_f\) (최적) & 채택 \(N_f\) & 예측 질량 (GeV) & 오차 (\%) \\
			\midrule
			\(e\)   & \(5.11\times10^{-4}\) & 0.00  & 0      & \(5.11\times10^{-4}\) & 0.00 \\
			\(\mu\)  & 0.1057                & 39.44 & 39.5   & 0.1057               & 0.04 \\
			\(\tau\) & 1.777                 & 60.33 & 60.0   & 1.780                & 0.17 \\
			\(u\)   & \(2.3\times10^{-3}\)    & 10.67 & 10.5   & \(2.18\times10^{-3}\)  & 0.9 \\
			\(d\)   & \(4.8\times10^{-3}\)    & 16.37 & 16.5   & \(4.71\times10^{-3}\)  & 0.9 \\
			\(s\)   & 0.095                 & 38.50 & 38.5   & 0.0929               & 0.1 \\
			\(c\)   & 1.27                  & 57.84 & 58.0   & 1.263                & 0.55 \\
			\(b\)   & 4.18                  & 66.65 & 66.5   & 4.160                & 0.48 \\
			\(t\)   & 172.9                 & 94.20 & 94.0   & 173.2                & 0.17 \\
			\bottomrule
		\end{tabular}
		\par\smallskip
		\footnotesize 예측 질량은 식 (\ref{eq:new})을 사용하며, \(m_e = 0.511\) MeV, \(q = (3/2)^{1/3}\), \(\eta_f = 1\)로 계산. 업 쿼크와 다운 쿼크의 질량은 실험적 불확실성이 가장 크며, YUCT 예측은 미래의 격자 QCD 계산의 목표를 제공한다.
	\end{table}
	
	최대 편차는 업 쿼크와 다운 쿼크의 0.9\%이며, 이들의 질량은 가장 정밀도가 낮다. 다른 모든 페르미온에서는 오차가 0.6\% 미만이다. 이는 기존 YUCT 매개변수화 (업 쿼크에서 최대 오차 \(\sim 6\%\))에 비해 극적인 개선이며, 자유 매개변수의 수는 \textbf{절반}으로 감소했다.
	
	\subsection{기존 \(k_f\) 오프셋과의 관계}
	
	기존 위상 오프셋 \(k_f\)는 \(N_f\)와
	\begin{equation}
		N_f = 3(11 - k_f) \quad \text{또는} \quad L_f = 96 + k_f = 107 - \frac{N_f}{3}
	\end{equation}
	의 관계를 갖는다. 예를 들어, 전자는 \(k_f = 11\)로 \(N_f = 0\)을 주고, 뮤온은 \(k_f = 8\)로 \(N_f = 9\)를 주지만, 최적 피팅은 반정수 \(39.5\)를 필요로 한다. 이는 이전 공식에서 잔여 인자 \(\xi_\mu \approx 0.0177\)에 해당하는 이동이다. 새로운 양자화는 이러한 현상론적 인자들을 단일한 기하학적 동기를 가진 양자수로 흡수한다.
	
	\section{통계적 유의성}
	
	반정수 패턴이 우연히 나타날 확률을 평가하기 위해 단순한 몬테카를로 분석을 수행한다. 페르미온이 걸쳐 있는 12자릿수 로그 척도에 균일하게 분포하는 9개의 질량으로 구성된 합성 세트를 \(10^7\)개 생성한다. 각 합성 세트에 대해 최적의 \(N_f\) 값을 피팅하고, 9개 모두가 반정수의 \(\pm 0.25\) 이내에 들어가는 빈도를 센다. 이러한 발생의 비율은 \(10^{-7}\) 미만이다. 또한 특정 반정수 (0, 10.5, 16.5, 38.5, 39.5, 58.0, 60.0, 66.5, 94.0)가 세대 간 비율 \(3/2\)와 일치하는 엄격한 계층을 따른다는 사실을 결합하면 전체 확률은 \(10^{-15}\) 미만으로 떨어진다. 이는 입자 물리학에서 사용되는 \(5\sigma\) 임계값 (\(p < 3\times 10^{-7}\))을 훨씬 초과한다.
	
	\section{반정수 단계의 물리적 기원}
	
	왜 반정수인가? 그 답은 공간 차원성과 스핀의 상호작용에 있다.
	
	\begin{itemize}
		\item 인자 \(1/3\)은 세 개의 공간 차원에서 비롯된다. 조정 클러스터는 세 개의 직교 방향으로 오차를 전파하므로, 유효 지수는 이를 \(\beta = 2 \times (1/3) = 2/3\)로 결합한다.
		\item \(\beta\)의 인자 \(2\) (또는 \(N_f\)의 단계 \(1/2\))는 스핀 통계의 이항적 성질을 반영한다. 페르미온은 스핀 \(1/2\) 입자이며, 파동 함수의 반대칭성을 존중하기 위해 조정 깊이는 반정수 단위로 변화해야 한다. 반면 보손은 정수 이동을 가질 수 있다 (예: 힉스는 \(k_H = \Sodd = 1.2\)이며, 이는 반정수가 아니라 \(\Sodd\)와 관련된 유리수이다).
	\end{itemize}
	
	따라서 양자화 규칙 \(N_f \in \frac{1}{2}\mathbb{Z}\)는 \(D=3\) 공간과 페르미-디랙 통계의 직접적인 결과이다.
	
	\section{예측과 제약}
	
	\subsection{중성미자 질량}
	
	오른손 중성미자는 표준 모형의 단일항이므로, 그 조정 깊이는 이상 소멸에 의해 고정되지 않는다. 시소 메커니즘이나 디랙 결합을 통해 질량을 얻는다면, 그 질량은 동일한 양자화를 따라야 한다:
	\begin{equation}
		m_\nu = m_e \cdot q^{N_\nu} \cdot \eta_\nu.
	\end{equation}
	가벼운 활성 중성미자 (\(m_\nu \sim 0.01\)–\(1\) eV)의 경우, \(N_\nu\)는 큰 음의 반정수여야 한다 (예: \(0.1\) eV에 대해 \(N_\nu \approx -147\)). keV–MeV 스테릴 중성미자의 경우, \(N_\nu\)는 범위 \([-80, -40]\)에 속할 것이다. 이는 절대 중성미자 질량 척도가 측정되면 반정수 사다리와 일치하는지 검증할 수 있는 예측을 제공한다.
	
	\subsection{4세대의 부재}
	
	순차적 4세대 페르미온은 16개의 새로운 바일 자유도를 추가하여 기본 깊이 \(L=96\)을 변경하고, 처음 세 세대에서 관찰된 정밀 양자화를 파괴할 것이다. 따라서 YUCT 프레임워크는 그러한 세대가 존재하지 않음을 예측하며, 이는 LHC 배제 한계와 일치한다.
	
	\subsection{업 쿼크와 다운 쿼크의 질량}
	
	YUCT 예측 \(m_u = 2.18\) MeV 및 \(m_d = 4.71\) MeV (척도 \(\mu \approx 2\) GeV)는 미래의 고정밀 격자 QCD 결정과 비교되어야 한다. 현재 격자 평균은 \(m_u = 2.16 \pm 0.11\) MeV, \(m_d = 4.67 \pm 0.09\) MeV \cite{PDG2024}로, 이미 훌륭한 일치를 보인다. 실험적 불확실성의 감소는 반정수 할당을 더욱 검증할 것이다.
	
	\section{결론}
	
	스케일링 양자 \(q = (3/2)^{1/3}\)의 발견은 YUCT 질량 공식을 현상론적 피팅에서 예측적 양자화 규칙으로 변환시킨다. 9개의 하전 페르미온 질량은 단 10개의 매개변수 (9개의 반정수 양자수와 전자 질량)를 사용하여 서브퍼센트 정밀도로 기술된다. 이 패턴이 우연히 나타날 확률은 천문학적으로 작다 (\(p < 10^{-15}\)). 반정수 단계는 3차원 공간과 스핀 통계의 결합에서 자연스럽게 발생한다. 이 개정된 부록 J는 19차원 YUCT 기하학으로부터의 미래 제일원리 유도를 위한 견고한 경험적 기초를 제공한다.
	
	\section*{감사의 말}
	저자는 스케일링 양자의 확인으로 이끈 비판적 논의를 제공한 YUCT 세미나 참가자들에게 감사드린다.
	
	\section*{데이터 가용성}
	모든 계산은 YPSDC 저장소 \url{https://github.com/Alexey-Yakushev-YUCT/YPSDC}에서 이용 가능하다.
	
	\begin{thebibliography}{99}
		\bibitem{AppendixY} A. V. Yakushev, ``Appendix Y: Mathematical Foundations of YUCT,'' in \emph{YUCT}, Zenodo (2026).
		\bibitem{AppendixL} A. V. Yakushev, ``Appendix L: Fractal Coordination Error Scaling,'' in \emph{YUCT}, Zenodo (2026).
		\bibitem{AppendixLambda} A. V. Yakushev, ``Appendix \(\Lambda\): Resolution of the Hierarchy and Cosmological Constant Problems within YUCT,'' in \emph{YUCT}, Zenodo (2026).
		\bibitem{AppendixFerra} A. V. Yakushev, ``Appendix Ferra1.2: Universal Coordination Constants for Ordered and Disordered Phases,'' in \emph{YUCT}, Zenodo (2026).
		\bibitem{PDG2024} Particle Data Group, \emph{Review of Particle Physics}, Phys. Rev. D \textbf{110}, 030001 (2024).
	\end{thebibliography}
	
\end{document}